% !TEX root = main.tex

\section{Appendix on weak-star limits}

\label{sec:app:weak star limits}

In this part, we collect useful statements that help the understanding of weak limit measures. Let $X$ be a locally compact Hausdorff topological space (see \cite[p31, p104]{lee2011IntroductionTopologicalManifolds}). We denote by $\BB$ the Borel $\sigma$-algebra on $X$, i.e., the $\sigma$-algebra generated by the open subsets of $X$ (see \cite[Section 15]{halmos1950MeasureTheory}). We recall the definition of a regular Borel measure.
\begin{definition}\cite[Section 52]{halmos1950MeasureTheory}
    A regular Borel measure $\mu : \BB \to [0, \infty]$ is a measure that is finite on compact sets, i.e., $\mu(C) < \infty$ for every compact set $C$, and that satisfies
        \begin{equation}
            \mu(A) = \inf \{\mu(U) : A \subseteq U, U \text{ is open}\} =  \sup \{\mu(U) : U \subseteq A, U \text{ is compact}\}.
        \end{equation}
\end{definition}
% For example, the measures $\nu_{\beta,m}$ defined in (\ref{eq:3:introduction to the formalism of Bernoulli convolution}) all are regular Borel measures. 
% We will show that the Bernoulli convolution $\nu_\beta$ defined by (\ref{eq:5:introduction to the formalism of Bernoulli convolution}) is also a regular Borel measure, because the weak limit preserves Borel regularity. Before stating the result, let us define the set of compactly supported continuous functions (Section 55 HALMOS)
% \begin{equation}
%     \CC_c(X) := \left\{ f: X \to \R : f \text{ is continuous, and } \overline{\{x : f(x) \neq 0\}}\text{ is compact} \right\}.
% \end{equation}
% Note that if $X$ is compact, $\CC_c(X)$ is simply the set of continuous functions.

Regular Borel measures are a central element of the following Lemmata, that are a key ingredient of the proof of Theorem \ref{thm:asymptotic behavior of card f x when the Bernoulli convolution is absolutely continuous}. Define the set of compactly supported continuous functions from $X$ to $\R$ as \cite[Section 55]{halmos1950MeasureTheory}
\begin{equation}
    \CC_c(X) := \left\{ f: X \to \R : f \text{ is continuous, and } \overline{\{x : f(x) \neq 0\}}\text{ is compact} \right\}.
\end{equation}
Note that if $X$ is compact, $\CC_c(X)$ is simply the set of continuous functions. 
\begin{lemma}
    Let $(\mu_n : \BB \to [0,\infty])_{n \in \N}$ be a sequence of regular Borel measures convering weakly, i.e., $\left(\int_X f d\mu_n\right)_{n \in \N}$ is converging in $\R$, for every $f \in \CC_c(X)$. Then, there exists a regular Borel measure $\mu : \BB \to [0,\infty]$ such that
    \begin{equation}
        \int_X f d\mu = \lim_{n\to \infty}\int_X f d\mu_n, \ \forall f \in \CC_c(X).
    \end{equation}
\end{lemma}
\begin{proof}
    A map $\Lambda : \CC_c(X) \to \R$ is said to be positive if $\Lambda(f) \geq 0$ when $f \geq 0$ and linear if $\Lambda(\alpha f + g) = \alpha \Lambda(f) + \Lambda(g)$ for all $\alpha \in \R$, $f,g \in \CC_c(X)$. For every measure $\mu : \BB \to [0,\infty]$, the functional defined by $\Lambda_\mu : f \mapsto \int_X f d\mu$ is obviously positive and linear. Define the functional $\Lambda : \CC_c(X) \to \R$ as 
    \begin{equation}
        \Lambda (f) := \lim_{n\to \infty}\int_X f d\mu_n, \ \forall f \in \CC_c(X).
    \end{equation}
    $\Lambda$ is well-defined since $(\mu_n)_{n \in \N}$ converges weakly. Moreover, $\Lambda$ is obviously positive and linear, by positivity and linearity of the limit. Further, there is a standard result, known as Riesz representation theorem, that states that for every positive linear functional $\Lambda : \CC_c(X) \to \R$, there exists a unique regular Borel measure $\mu$ such that $\Lambda = \Lambda_\mu$ (see  \cite[Section 56, Theorems D and E]{halmos1950MeasureTheory}). This concludes the proof.
\end{proof}

\begin{lemma}
    Suppose that $X$ is compact. Let $(\nu_n)_{n \in \N}$ be a sequence of regular Borel measures converging weakly to a Borel regular measure $\nu$, i.e.,
    \begin{equation}
        \int_Xfd\mu = \lim_{n \to \infty} \int_X fd\mu_n, \ \forall f \in \CC_c(X).
    \end{equation}
    Then, for every compact set $C$, we have
    \begin{equation}
        \limsup_{n \to \infty} \mu_n(C) \leq \mu(C).
    \end{equation}
\end{lemma}
\begin{proof} 
    Let $(\nu_n)_{n \in \N}$ be a sequence of regular Borel measures converging weakly to the Borel regular measure $\nu$ and $C$ be a closed set. Let $U$ be an open set, so that $C \subseteq U$. Since $U$ is open, then $X \backslash U$ is closed and disjoint from $C$. Since $X$ is compact and Hausdorff, then it is normal (see \cite[p111]{lee2011IntroductionTopologicalManifolds} and \cite[Exercise 4.78]{lee2011IntroductionTopologicalManifolds}), so we can apply Urysohn's Lemma \cite[Lemma 4.82]{lee2011IntroductionTopologicalManifolds}: there exists a continuous function $f_U : X \to [0,1]$ such that $f(C) = \{1\}$ and $f(X\backslash U) = \{0\}$. Note that since 
    \begin{equation}
        \chi_C(x) \leq f_U(x) \leq \chi_U(x), \ \forall x \in X,
    \end{equation}
    where $\chi_A$ denotes the characteristic function of $A \in \BB$, then
    \begin{equation}
        \mu_n(C) = \int_X \chi_C d\mu_n \leq \int_Xf_Ud\mu_n,
    \end{equation}
    and therefore
    \begin{equation}
        \limsup_{n \to \infty}\mu_n(C) \leq \limsup_{n \to \infty}\int_Xf_Ud\mu_n = \int_Xf_Ud\mu \leq \int_X\chi_U d\mu = \mu(U).
    \end{equation}
    Since the left-hand side of the inequality does not depend on $U$, we get
    \begin{equation}
        \limsup_{n \to \infty}\mu_n(C) \leq \inf_{\substack{U \text{ open},\\ U \subseteq C}} \mu(U) = \mu(C),
    \end{equation}
    where the last equality follows by Borel regularity of $\mu$.
\end{proof}

Since $\R$ is a locally compact Hausdorff topological space, and every closed interval $[a,b]$, $a,b \in \R$ is a compact subset of $\R$, we can use the above Lemma to prove Theorem \ref{thm:asymptotic behavior of card f x when the Bernoulli convolution is absolutely continuous}. 